% A1 ----------------------------------------------------------
\begin{frame}[t]{Replay protocol and information boundaries}
\vspace{0.3cm}
\begin{columns}[T,onlytextwidth]
\column{0.49\textwidth}
\keyterm{Data separation}
{\small
\begin{itemize}\setlength\itemsep{0.45em}
\item Group duplicates and common source items before splitting.
\item Separate reference, development, calibration, and test pools.
\item Fit topic groups on development data only.
\item Freeze horizon, window, stride, and replay design before testing.
\end{itemize}}
\column{0.48\textwidth}
\keyterm{Interpretation}
{\small
\begin{itemize}\setlength\itemsep{0.45em}
\item Reordering a fixed corpus adds no independent queries.
\item Calibrate the full alarm and update policy.
\item Verified updates require an explicit information source.
\item Query-only monitoring needs assumptions about benign change.
\end{itemize}}
\end{columns}
\par\vspace{0.4cm}
\smallnote{Identical observable query distributions cannot reveal different intentions from query content alone.}
\sourceline{Proposed protocol. Calibrate the complete alarm policy before held-out evaluation.}
\end{frame}

% A2 ----------------------------------------------------------
\begin{frame}[t]{Related work and the proposed increment}
\vspace{0.2cm}
\renewcommand{\arraystretch}{1.4}
{\small
\begin{tabularx}{\textwidth}{@{}L{3.4cm}L{5.0cm}Y@{}}
\toprule
\keyterm{Existing work} & \keyterm{What it contributes} & \keyterm{Role in the proposed study}\\
\midrule
Liu, Guo \& Dong\newline 2026 &
MMD on query windows with benign calibration. &
Released starting point.\\
Flynn \& Yoo\newline CDC 2019 &
KCUSUM: kernel sequential change detection. &
Accumulation baseline.\\
Cobb et al.\newline AISTATS 2022 &
CALM: sequential false-detection calibration. &
Calibration baseline.\\
\bottomrule
\end{tabularx}\par}
\vspace{0.45cm}
{\small Proposed increment: continuous API replay, followed by any unresolved reference-update failure under explicit information assumptions.\par}
\sourceline{
\href{https://arxiv.org/abs/2606.05725}{MMD: arXiv 2606.05725}\quad
\href{https://arxiv.org/abs/1903.01661}{KCUSUM: arXiv 1903.01661}\quad
\href{https://proceedings.mlr.press/v151/cobb22a.html}{CALM: PMLR 151, 226--239}.}
\end{frame}

% A3 ----------------------------------------------------------
\begin{frame}[t]{Provenance and applicability of reference updates}
\vspace{0.1cm}
\begin{columns}[T,onlytextwidth]
\column{0.48\textwidth}
\keyterm{LangSkills: traceable artifacts}
{\small
\begin{itemize}\setlength\itemsep{0.5em}
\item Records source references and generation metadata.
\item Validates skill structure and metadata.
\item These checks do not establish scientific validity.
\end{itemize}}
\column{0.48\textwidth}
\keyterm{HazardWeaver: explicit conditions}
{\small
\begin{itemize}\setlength\itemsep{0.5em}
\item Links applicability conditions to evidence.
\item Distinguishes satisfied, violated, and unresolved conditions.
\item Unresolved required evidence blocks immediate use.
\end{itemize}}
\end{columns}
\par\vspace{0.2cm}
{\small\keyterm{Application to this proposal:} record reference provenance, verification evidence, and the benign-change regime supported by calibration.\par}
\vspace{0.1cm}
\smallnote{This documentation principle is proposed here. Neither source establishes its effectiveness for API extraction monitoring.}
\sourceline{\href{https://github.com/LabRAI/LangSkills}{LangSkills} (source artifacts and validator).\quad
\href{https://arxiv.org/html/2610.03591v1}{Zhu et al., HazardWeaver, 2026} (Sec.\ 3.1 and Appendix J).}
\end{frame}
